\section{真题代练}

\subsection{使用说明}
\begin{enumerate}
  \item 每题先\textbf{独立完成}（选择题 2 分钟，非选择题 8--12 分钟），再对照解析
  \item 重点关注\textbf{思路分析}和\textbf{答题规范}
  \item 标记错题，定期重做
\end{enumerate}

\subsection{一、分子与细胞}

\textbf{例 1}（2024 新课标Ⅰ卷）下列关于细胞中化合物的叙述，错误的是
\begin{enumerate}
  \item[A.] 磷脂是所有细胞必不可少的脂质
  \item[B.] 组成淀粉和糖原的单体都是葡萄糖
  \item[C.] 蛋白质变性后仍能与双缩脲试剂反应
  \item[D.] DNA 分子中 G--C 碱基对越多，热稳定性越弱
\end{enumerate}

\textbf{思路分析}：逐项判断——A 磷脂是膜结构必需成分（正确）；B 淀粉糖原均为多糖，单体葡萄糖（正确）；C 变性不破坏肽键（正确）；D G--C 间三个氢键比 A--T 更稳定（错误）。

\textbf{解析}：D 错误——G--C 之间形成 3 个氢键，A--T 之间形成 2 个氢键，G--C 含量越高 DNA 热稳定性越强（熔点越高）。A 正确——所有细胞都有膜结构，膜的基本骨架是磷脂双分子层。B 正确——淀粉和糖原均由葡萄糖通过糖苷键连接而成。C 正确——双缩脲试剂与肽键反应，变性不破坏肽键。

\textbf{答案}：\boxed{D}

\textbf{例 2}（2023 新课标Ⅰ卷·节选）将某植物叶片置于密闭装置中，测定不同光照强度下 CO$_2$ 吸收速率，结果如下表：

\begin{center}
\begin{tabular}{c|c|c|c|c|c|c}
光照强度（klx） & 0 & 2 & 4 & 6 & 8 & 10 \\ \hline
CO$_2$ 吸收速率（mg/h） & -4 & 0 & 4 & 8 & 10 & 10
\end{tabular}
\end{center}

（1）光照强度为 0 时，CO$_2$ 吸收速率为负值的原因是\_\_\_\_\_\_。
（2）光补偿点为\_\_\_\_\_\_klx，此时该叶片的光合速率\_\_\_\_\_\_（填"$>$""$=$"或"$<$"）呼吸速率。
（3）光照强度为 8 klx 时，限制光合速率的主要因素是\_\_\_\_\_\_。

\textbf{解析}：
（1）光照强度为 0 时无光合作用，只进行细胞呼吸释放 CO$_2$，故 CO$_2$ 吸收速率为负值
（2）光补偿点是光合速率等于呼吸速率时的光照强度，表中 CO$_2$ 吸收速率为 0 时对应 2 klx，此时光合速率 = 呼吸速率
（3）8 klx 时已达光饱和点（10 klx 与 8 klx 速率相同），限制因素转为 CO$_2$ 浓度或温度等非光因素

\textbf{答案}：（1）呼吸作用释放 CO$_2$ 多于光合作用吸收 （2）2，= （3）CO$_2$ 浓度

\textbf{例 3}（2022 新课标Ⅰ卷）下列关于细胞呼吸的叙述，正确的是
\begin{enumerate}
  \item[A.] 有氧呼吸第一阶段产生 CO$_2$
  \item[B.] 无氧呼吸第二阶段产生少量 ATP
  \item[C.] 线粒体是细胞呼吸的唯一场所
  \item[D.] 有氧呼吸第三阶段产生大量 ATP
\end{enumerate}

\textbf{解析}：A 错误——有氧呼吸第一阶段产生丙酮酸和[H]，不产生 CO$_2$（CO$_2$ 在第二阶段产生）；B 错误——无氧呼吸第二阶段不产生 ATP；C 错误——有氧呼吸第一阶段在细胞质基质进行；D 正确——第三阶段[H]与 O$_2$ 结合生成水，释放大量能量。

\textbf{答案}：\boxed{D}

\subsection{二、遗传与进化}

\textbf{例 4}（2024 新课标Ⅰ卷）某种植物的花色由两对独立遗传的等位基因 A/a 和 B/b 控制，A 和 B 同时存在时开红花，否则开白花。现有纯合红花植株与纯合白花植株杂交，F$_1$ 全为红花，F$_1$ 自交得 F$_2$。请回答：
（1）亲本的基因型为\_\_\_\_\_\_和\_\_\_\_\_\_。
（2）F$_2$ 中红花:白花的比例为\_\_\_\_\_\_。
（3）F$_2$ 白花中纯合子的比例为\_\_\_\_\_\_。

\textbf{解析}：
（1）纯合红花（AABB）$\times$ 纯合白花（因 A 和 B 同时存在才为红花，白花纯合可能为 AAbb、aaBB、aabb，但 F$_1$ 全为红花，故亲本白花应为 aabb）
（2）F$_1$ AaBb 自交：A$\_$B$\_$（红花）: A$\_$bb : aaB$\_$ : aabb = 9:3:3:1，故红花:白花 = 9:7
（3）F$_2$ 白花（7 份）中纯合子包括 AAbb（1）、aaBB（1）、aabb（1），共 3 份，故比例为 3/7

\textbf{答案}：（1）AABB $\times$ aabb （2）9:7 （3）3/7

\textbf{例 5}（2023 新课标Ⅰ卷）下列关于人类遗传病的叙述，正确的是
\begin{enumerate}
  \item[A.] 单基因遗传病是由一个致病基因引起的
  \item[B.] 21 三体综合征属于染色体数目异常遗传病
  \item[C.] 通过遗传咨询可检测胎儿是否患遗传病
  \item[D.] 近亲结婚会增加常染色体显性遗传病的发病率
\end{enumerate}

\textbf{解析}：A 错误——单基因遗传病是由一对等位基因引起的；B 正确——21 三体综合征为额外多一条 21 号染色体；C 错误——遗传咨询不能直接检测，产前诊断（羊水检查等）才可检测；D 错误——近亲结婚主要增加隐性遗传病的发病率。

\textbf{答案}：\boxed{B}

\textbf{例 6}（2021 新课标Ⅰ卷）果蝇的灰身（B）对黑身（b）为显性，长翅（V）对残翅（v）为显性，两对基因均位于常染色体上。现有灰身长翅果蝇与黑身残翅果蝇杂交，F$_1$ 中灰身长翅:灰身残翅:黑身长翅:黑身残翅 = 1:1:1:1。请分析：
（1）该实验结果说明\_\_\_\_\_\_。
（2）若 F$_1$ 中灰身长翅果蝇与黑身残翅果蝇杂交，子代出现灰身长翅:黑身残翅 = 1:1，其原因最可能是\_\_\_\_\_\_。

\textbf{解析}：
（1）测交结果 1:1:1:1，说明灰身长翅亲本产生了四种比例相等的配子，即两对基因位于非同源染色体上（自由组合）或连锁但交换率足够高
（2）后代表现型比例 1:1（亲本型），说明灰身长翅果蝇只产生了两种配子（BV 和 bv），即两对基因完全连锁

\textbf{答案}：（1）两对基因遵循自由组合定律（2）两对基因位于同一对同源染色体上，完全连锁

\subsection{三、稳态与调节}

\textbf{例 7}（2024 新课标Ⅰ卷）下列关于神经调节的叙述，错误的是
\begin{enumerate}
  \item[A.] 静息电位时 K$^+$ 外流属于协助扩散
  \item[B.] 兴奋在反射弧中单向传递
  \item[C.] 突触前膜释放神经递质需消耗 ATP
  \item[D.] 突触后膜上的受体与递质结合后一定引起 Na$^+$ 内流
\end{enumerate}

\textbf{解析}：A 正确——K$^+$ 通过钾离子通道顺浓度梯度外流（协助扩散）；B 正确——反射弧中存在突触，突触传递是单向的；C 正确——突触小泡与突触前膜融合释放递质需能量；D 错误——抑制性递质引起 Cl$^-$ 内流（使膜电位更负），而非 Na$^+$ 内流。

\textbf{答案}：\boxed{D}

\textbf{例 8}（2023 新课标Ⅰ卷·节选）人体在寒冷环境中，甲状腺激素分泌增加。请回答：
（1）寒冷刺激下，\_\_\_\_\_\_（器官）分泌促甲状腺激素释放激素（TRH）增加，作用于垂体，促进\_\_\_\_\_\_的分泌。
（2）甲状腺激素的作用包括\_\_\_\_\_\_、\_\_\_\_\_\_。
（3）当甲状腺激素含量过高时，会\_\_\_\_\_\_下丘脑和垂体的分泌活动，这属于\_\_\_\_\_\_调节。

\textbf{解析}：
（1）寒冷刺激下丘脑体温调节中枢 $\to$ 下丘脑分泌 TRH $\to$ 垂体分泌 TSH
（2）甲状腺激素主要功能：促进细胞代谢（物质氧化分解产热）、促进生长发育（尤其神经系统）、提高神经兴奋性
（3）甲状腺激素过多时反馈抑制下丘脑和垂体分泌，属于负反馈调节

\textbf{答案}：（1）下丘脑，促甲状腺激素（TSH） （2）促进新陈代谢（产热），促进生长发育 （3）抑制，负反馈

\textbf{例 9}（2022 新课标Ⅰ卷）下列关于免疫的叙述，正确的是
\begin{enumerate}
  \item[A.] 浆细胞可以识别抗原
  \item[B.] 记忆细胞不能直接产生抗体
  \item[C.] 辅助 T 细胞只参与细胞免疫
  \item[D.] 过敏反应是免疫防御功能过弱的表现
\end{enumerate}

\textbf{解析}：A 错误——浆细胞不能识别抗原，只能产生抗体；B 正确——记忆细胞受抗原刺激后增殖分化为浆细胞，浆细胞产生抗体；C 错误——辅助 T 细胞既参与体液免疫又参与细胞免疫；D 错误——过敏反应是免疫防御功能过强。

\textbf{答案}：\boxed{B}

\subsection{四、生物与环境}

\textbf{例 10}（2024 新课标Ⅰ卷）下列关于种群和群落的叙述，错误的是
\begin{enumerate}
  \item[A.] 年龄结构可预测种群数量的变化趋势
  \item[B.] 群落的垂直结构提高了群落利用环境资源的能力
  \item[C.] 初生演替和次生演替的起始条件不同
  \item[D.] 种间竞争总是导致一方灭绝
\end{enumerate}

\textbf{解析}：A 正确——年龄结构（增长型/稳定型/衰退型）反映种群出生率与死亡率关系；B 正确——垂直分层使不同生物利用不同层次资源；C 正确——初生演替起点无土壤条件，次生演替保留土壤；D 错误——种间竞争可以共存（生态位分化），不一定导致一方灭绝。

\textbf{答案}：\boxed{D}

\textbf{例 11}（2023 新课标Ⅰ卷·节选）某生态系统中存在食物链：草 $\to$ 蝗虫 $\to$ 青蛙 $\to$ 蛇。回答：
（1）该食物链中，蝗虫属于第\_\_\_\_\_\_营养级，属于\_\_\_\_\_\_（成分）。
（2）若青蛙的能量 1/2 来自蝗虫，1/2 来自蜘蛛（草$\to$蝗虫$\to$蜘蛛$\to$青蛙），则青蛙增重 1 kg，最多需要草\_\_\_\_\_\_kg。
（3）能量沿食物链流动的特点是\_\_\_\_\_\_。

\textbf{解析}：
（1）蝗虫以草（第一营养级）为食，为第二营养级，属于消费者（初级消费者）
（2）最多需要草——按 10\% 传递效率计算。青蛙增重 1 kg：0.5 kg 来自蝗虫 $\to$ 需蝗虫 5 kg $\to$ 需草 50 kg；0.5 kg 来自蜘蛛 $\to$ 需蜘蛛 5 kg $\to$ 需蝗虫 50 kg $\to$ 需草 500 kg。共需草 50 + 500 = 550 kg
（3）单向流动、逐级递减

\textbf{答案}：（1）二，消费者 （2）550 （3）单向流动、逐级递减

\subsection{五、生物技术与工程}

\textbf{例 12}（2024 新课标Ⅰ卷）下列关于微生物培养的叙述，正确的是
\begin{enumerate}
  \item[A.] 培养基中必须添加琼脂作为凝固剂
  \item[B.] 接种环和涂布器均采用灼烧灭菌
  \item[C.] 稀释涂布平板法可用于微生物计数
  \item[D.] 选择培养基只能筛选目标微生物
\end{enumerate}

\textbf{解析}：A 错误——液体培养基不需要琼脂；B 正确——接种环和涂布器在使用前后均需灼烧灭菌；C 正确——稀释涂布平板法通过菌落数计算活菌数量；D 错误——选择培养基可筛选目标微生物（生长）也可抑制非目标微生物。

\textbf{答案}：\boxed{C}

\textbf{例 13}（2023 新课标Ⅰ卷·节选）利用基因工程技术将人胰岛素基因导入大肠杆菌，实现胰岛素的大规模生产。回答：
（1）获取人胰岛素基因的方法有\_\_\_\_\_\_（答 2 种）。
（2）构建基因表达载体时，需在目的基因前添加\_\_\_\_\_\_以便于转录。
（3）常用\_\_\_\_\_\_处理大肠杆菌，使其成为感受态细胞，以利于吸收重组 DNA 分子。
（4）如何从培养液中分离纯化胰岛素？\_\_\_\_\_\_。

\textbf{解析}：
（1）获取目的基因方法：从 cDNA 文库中获取（以 mRNA 为模板逆转录）、PCR 扩增、化学人工合成
（2）启动子是 RNA 聚合酶结合位点，驱动基因转录
（3）Ca$^{2+}$（氯化钙）处理增加细胞膜通透性
（4）胰岛素分泌到培养液中，可通过盐析、凝胶色谱、亲和层析等方法纯化

\textbf{答案}：（1）从 cDNA 文库中获取，PCR 扩增 （2）启动子 （3）Ca$^{2+}$ 溶液 （4）盐析/层析

\subsection{六、综合实验与探究}

\textbf{例 14}（2024 新课标Ⅰ卷）某兴趣小组探究不同浓度 2,4-D（生长素类似物）对某植物插条生根的影响。请回答：
（1）实验中用浸泡法处理插条时，处理时间\_\_\_\_\_\_（填"属于"或"不属于"）无关变量。
（2）预实验的目的是\_\_\_\_\_\_。
（3）正式实验中，各组应控制\_\_\_\_\_\_等无关变量相同（答 2 点）。
（4）若实验结果出现低浓度促进生根、高浓度抑制生根的现象，说明 2,4-D 的作用具有\_\_\_\_\_\_。

\textbf{解析}：
（1）处理时间是无关变量，各组应保持一致
（2）预实验为正式实验摸索条件（浓度范围），检验实验设计的科学性和可行性
（3）无关变量包括：插条的生理状态（粗细、年龄）、培养温度、光照条件、处理时间等
（4）低浓度促进、高浓度抑制——两重性

\textbf{答案}：（1）属于 （2）探索最适浓度范围（避免盲目性） （3）插条生理状态、培养条件 （4）两重性

\textbf{例 15} 为验证"甲状腺激素能促进动物个体发育"，请以蝌蚪为实验材料设计实验方案。

\textbf{方案}：
\begin{enumerate}
  \item 取生理状态相同的蝌蚪若干只，随机均分为三组：A 组、B 组、C 组
  \item A 组（实验组）：饲喂含甲状腺激素的饲料；B 组（对照组）：饲喂等量不含甲状腺激素的普通饲料；C 组（抑制组）：饲喂含甲状腺激素合成抑制剂的饲料
  \item 在相同且适宜的条件下培养，每天观察蝌蚪的变态发育情况（记录四肢出现时间、尾的消失时间等）
  \item 结果预期：A 组蝌蚪发育最快（提前变态），B 组正常发育，C 组发育迟缓或不发育
  \item 结论：甲状腺激素能促进蝌蚪个体发育
\end{enumerate}
